% =============================================================================
% TEMPLATE DE LISTA / SIMULADO  —  estilo Avaliação Final (MPE)
% Preencha o conteúdo; a estrutura e a regra de compilação já estão prontas.
%
% DOIS MODOS (um único arquivo-fonte):
%   \solucoesfalse  -> gera a LISTA em branco (espaço para resolver)
%   \solucoestrue   -> gera o GABARITO + RUBRICA (mesmo arquivo)
%
% Compilar (PDF copiável, regra global lmodern+cmap):
%   pdflatex lista.tex   (rode 2x)   ; verifique com: pdffonts lista.pdf  (tudo Type 1)
% =============================================================================
\documentclass[11pt,a4paper]{article}
\usepackage[utf8]{inputenc}
\usepackage[T1]{fontenc}
\usepackage{lmodern}   % fontes vetoriais Type 1 (substitui bitmap PK)
\usepackage{cmap}      % ToUnicode CMap -> texto e acentos copiaveis
\usepackage[brazil]{babel}
\usepackage{amsmath,amssymb,amsfonts}
\usepackage[a4paper,margin=2.2cm]{geometry}
\usepackage{enumitem}
\usepackage{fancyhdr}
\usepackage{comment}
\usepackage{xcolor}
\usepackage{microtype}

% ====================== CONFIGURAÇÃO ======================
\newif\ifsolucoes
\solucoesfalse                 % <<<< \solucoestrue para GABARITO; \solucoesfalse para LISTA

\newcommand{\disciplina}{Microeconomia I --- MPE}
\newcommand{\listatema}{\textit{<tema da lista>}}
\newcommand{\espacoresol}{4.5cm}   % altura do espaço de resolução na lista em branco

% ====================== MAQUINARIA (não precisa editar) ======================
\ifsolucoes\includecomment{sol}\else\excludecomment{sol}\fi

\pagestyle{fancy}\fancyhf{}
\lhead{\footnotesize\disciplina}
\rhead{\footnotesize Lista\ifsolucoes{} --- Soluções e Rubrica\fi}
\cfoot{\footnotesize\thepage}
\renewcommand{\headrulewidth}{0.3pt}
\setlength{\headheight}{14pt}

\setlist{topsep=2pt,itemsep=2pt}
\newcounter{questao}

% Cabeçalho de bloco temático
\newcommand{\bloco}[3]{\bigskip\noindent
  {\large\bfseries\color{blue!55!black}Bloco #1 --- #2}\hfill\textit{(#3)}\par
  \rule{\linewidth}{0.3pt}\smallskip}

% Questão (numeração automática)
\newcommand{\questao}[1]{\refstepcounter{questao}\medskip\noindent
  \textbf{Questão \thequestao.}~#1\par}

% Item objetivo de múltipla escolha: \mc{enunciado}{lista de \item}
\newcommand{\mc}[1]{#1\par}
\newenvironment{alternativas}{\begin{enumerate}[label=\Alph*.,leftmargin=2.2em,topsep=1pt]}{\end{enumerate}}

% Gabarito objetivo (só aparece no modo soluções)
\newcommand{\gabarito}[1]{\ifsolucoes\par\noindent\textbf{Gabarito:}~#1\fi}

% Espaço de resolução na lista em branco
\newcommand{\resolespaco}{\ifsolucoes\else\par\textit{Resolução:}\vspace{\espacoresol}\par\fi}

% Passo de resolução (use dentro de sol)
\newcommand{\passo}[2]{\par\smallskip\noindent\textbf{Passo --- #1.}~#2}
\newcommand{\intuicao}[1]{\par\smallskip\noindent\textit{Intuição econômica.}~#1}
\newcommand{\refbib}[1]{\par\smallskip\noindent{\footnotesize\color{black!60}\textbf{Ref:} #1}}
% Rubrica de correção
\newcommand{\rubrica}[1]{\par\smallskip\noindent\textbf{Rubrica de correção.}\begin{itemize}[leftmargin=1.4em,topsep=1pt]#1\end{itemize}}

% =============================================================================
\begin{document}
\begin{center}
  {\Large\bfseries \disciplina}\\[2pt]
  {\large Lista de Exercícios: \listatema}\ifsolucoes\\[2pt]\textit{Soluções e Rubrica}\fi
\end{center}
\vspace{2pt}

{\footnotesize\noindent\textbf{Instruções.} Itens de múltipla escolha: assinale a única alternativa correta
(sem penalização por erro). Itens \emph{Verdadeiro/Falso}: a cada 2 itens errados anula-se 1 correto
(nota $=\max\{0,\text{acertos}-\text{erros}/2\}\times$ valor); itens em branco não contam.
Questões discursivas: mostre os passos.\par}

% =================== EXEMPLO (substitua pelo conteúdo gerado) ===================
\bloco{1}{Teoria do Consumidor}{Aulas 1--3 \,\textbullet\, 25 pontos}

\questao{(10 pontos) Itens objetivos.}

\textbf{a)} (3 pontos) A Petrobras reajusta o preço da gasolina ($x_1$). Um consumidor com
$u(x_1,x_2)=x_1^{1/2}x_2^{1/2}$, renda $m=120$, preços $p_1=2,\ p_2=1$. A demanda
marshalliana por gasolina $x_1^M$ vale:
\begin{alternativas}
\item $20$
\item $30$
\item $40$
\item $60$
\end{alternativas}
\gabarito{B (\,$x_1^M=\tfrac{\alpha m}{p_1}=\tfrac{0{,}5\cdot120}{2}=30$\,)}
\begin{sol}
\passo{demanda CD}{Para $u=x_1^\alpha x_2^{1-\alpha}$ com $\alpha=\tfrac12$, $x_1^M=\dfrac{\alpha m}{p_1}=\dfrac{0{,}5\cdot 120}{2}=30$.}
\refbib{N\&S 12e, Cap.~4 \S4.1; ZaE Cap.~4.}
\end{sol}

\medskip
\textbf{c)} (2 pontos, V ou F) Se a gasolina é um bem de Giffen para esse consumidor, então ela é
necessariamente um bem inferior.
\gabarito{V}
\begin{sol}
\passo{critério}{Pela equação de Slutsky, $\partial x^M/\partial p>0$ exige efeito-renda negativo
dominante, o que requer bem inferior ($\partial x/\partial m<0$). Logo todo Giffen é inferior (recíproca falsa).}
\refbib{N\&S 12e, Cap.~5 \S5.1; MWG \S3.G.}
\end{sol}

\questao{(15 pontos) Pigou. A CETESB regula uma fábrica com demanda inversa
$P(q)=24-2q$, custo marginal privado $\mathrm{MPC}=6$ e custo externo marginal $\mathrm{MEC}(q)=q$.}

\textbf{a)} (5 pontos) Determine a quantidade socialmente ótima $q^*$ e o imposto de Pigou $t^*$.
\resolespaco
\begin{sol}
\passo{ótimo social}{Iguala benefício marginal ao custo marginal social:
$24-2q = 6 + q \Rightarrow q^*=6$.}
\passo{imposto}{$t^*=\mathrm{MEC}(q^*)=q^*=6$.}
\intuicao{O imposto internaliza o dano avaliado \emph{no ótimo social}, não no privado.}
\refbib{N\&S 12e, Cap.~19 \S19.2; MWG \S11.B.}
\rubrica{
\item Iguala BMg a CMg social ($24-2q=6+q$): 2 pts.
\item Resolve $q^*=6$: 1 pt.
\item $t^*=\mathrm{MEC}(q^*)=6$: 2 pts.
\item Avaliar MEC no nível privado em vez do ótimo: máx.\ 2 pts no item (teto conceitual).
}
\end{sol}

\textbf{b)} (5 pontos) Calcule o peso-morto da inação (produzir $q^{priv}$ em vez de $q^*$).
\resolespaco
\begin{sol}
\passo{nível privado}{Sem imposto: $24-2q=6 \Rightarrow q^{priv}=9$.}
\passo{DWL}{$\mathrm{DWL}=\int_{q^*}^{q^{priv}}[\mathrm{MEC}(q)-(P(q)-\mathrm{MPC})]\,dq
=\int_{6}^{9}[q-(18-2q)]\,dq=\int_6^9(3q-18)\,dq=13{,}5$.}
\refbib{N\&S 12e, Cap.~19 \S19.2.}
\rubrica{
\item Encontra $q^{priv}=9$: 1 pt.
\item Monta a área entre dano marginal e lucro marginal: 2 pts.
\item Integra corretamente ($=13{,}5$): 2 pts; erro aritmético com método certo: $-1$ pt.
}
\end{sol}

\end{document}
